\documentclass[12pt, a4paper]{article}

% ----- GÓI NGÔN NGỮ VÀ FONT TIẾNG VIỆT -----
\usepackage[utf8]{inputenc}
\usepackage[T5]{fontenc}
\usepackage[vietnamese]{babel}

% ----- CÁC GÓI TOÁN HỌC -----
\usepackage{amsmath}
\usepackage{amssymb}
\usepackage{amsfonts}

% ----- CÁC GÓI BỐ CỤC VÀ KHÁC -----
\usepackage[margin=1in]{geometry} % Căn lề
\usepackage{graphicx}             % Chèn hình ảnh
\usepackage{array}                % Cải tiến môi trường tabular
\usepackage{parskip}              % Thụt lề đoạn = 0, thêm khoảng cách giữa các đoạn
\usepackage{fancyhdr}             % Tùy chỉnh đầu trang/chân trang

%

% ----- CÀI ĐẶT ĐÁNH SỐ TRANG -----
\pagestyle{fancy}
\fancyhf{} % Xóa hết header/footer
\fancyfoot[C]{\thepage} % Đặt số trang ở giữa footer
\renewcommand{\headrulewidth}{0pt} % Bỏ đường kẻ ở header

% ----- THÔNG TIN TÀI LIỆU -----
%\title{BẢO TOÀN MOMEN ĐỘNG LƯỢNG CỦA CHẤT ĐIỂM}
\author{}
\date{} % Không hiển thị ngày

\begin{document}
	
	% --- NỘI DUNG VĂN BẢN ---
	% \noindent được sử dụng để xóa thụt đầu dòng của đoạn văn
	\noindent 
	% \textbf dùng để in đậm số thứ tự của bài
	\section*{Bài 1}

Một quả cầu nhỏ trượt không ma sát theo bề mặt bên trong của một phễu hình nón thẳng đứng. 
Tại thời điểm ban đầu, quả cầu ở độ cao $h_0$ và vận tốc nằm ngang $v_0$. 
Tìm $v_0$, nếu biết rằng, khi tiếp tục chuyển động, quả cầu lên đến độ cao $h$, sau đó bắt đầu hạ xuống. 
Tìm vận tốc $v$ của quả cầu tại điểm cao nhất.
\begin{figure}[h]
	\centering
	\includegraphics[width=0.3\linewidth]{"../../../itrithuc/OneDrive - Trong/OneDrive/0 AI LATEX/THU VIEN ANH/screenshot008"}
	\caption{}
	\label{fig:screenshot008}
\end{figure}
\section*{Bài 2}

Một bán cầu rỗng bán kính $R$, mặt trong nhẵn, được giữ cố định trên mặt đất sao cho mặt hở hướng lên trên. Một vật nhỏ ở điểm cao nhất của mặt trong bán cầu và được truyền một vận tốc $v_0$ theo phương tiếp tuyến với miệng bán cầu (hình 1). Cho gia tốc trọng trường là $g$.

\begin{enumerate}
	\item[a.] Tìm độ cao nhỏ nhất của quỹ đạo chuyển động của vật so với mặt đất.
	\item[b.] Tìm vận tốc lớn nhất của vật nhỏ trong quá trình chuyển động.
\end{enumerate}
\begin{figure}[h!]
	\centering
	\includegraphics[width=0.3\linewidth]{"../../../itrithuc/OneDrive - Trong/OneDrive/0 AI LATEX/THU VIEN ANH/screenshot002"}
	\caption{}
	\label{fig:screenshot002}
\end{figure}

\subsection*{Đáp số}

a/ $h_{\text{min}} = R(1-\cos\theta) = R - \frac{v_0^2}{4g} \left( \sqrt{\frac{16g^2R^2}{v_0^4} + 1} - 1 \right)$

\vspace{0.5cm} % Thêm một chút khoảng cách

b/ $v = \sqrt{\frac{v_0^2 + \sqrt{16g^2R^2 + v_0^4}}{2}}$
\section*{Bài 3.} Trên một mặt phẳng ngang nhẵn có một vật nhỏ khối lượng $m$
	chuyển động ; vật này được buộc vào đầu một sợi dây không dãn ;
	đầu kia của dây được kéo qua lỗ O (hình vẽ) với vận tốc
	kéo không đổi. Xác định sức căng của dây theo khoảng cách $r$
	giữa vật và lỗ O, nếu với $r = r_0$, vận tốc góc của dây bằng $\omega_0$.
	% TODO: \usepackage{graphicx} required
	
	\begin{figure}[h]
		\centering
		\includegraphics[width=0.5\linewidth]{"../../../itrithuc/OneDrive - Trong/OneDrive/0 AI LATEX/THU VIEN ANH/screenshot006"}
		\caption{}
		\label{fig:screenshot006}
	\end{figure}
	

	\section*{Bài 4}
	
	Một vật khối lượng $m$ được cột chặt vào đầu của 1 sợi dây. 
	Vật dịch chuyển không ma sát trên bàn, và dây được xuyên qua 1 lỗ nhỏ trên bàn (bỏ qua ma sát) như hình vẽ, dưới đó có 1 con chó kéo dây làm cho dây luôn bị căng.
	
	Ban đầu vật chuyển động theo quỹ đạo tròn, với động năng $W_d$. 
	Dây được kéo từ từ, cho đến khi bán kính của vòng giảm đi một nửa. 
	Tìm công mà con chó đã thực hiện.
	\begin{figure}[h]
		\centering
		\includegraphics[width=0.5\linewidth]{"../../../itrithuc/OneDrive - Trong/OneDrive/0 AI LATEX/THU VIEN ANH/screenshot007"}
		\caption{}
		\label{fig:screenshot007}
	\end{figure}
\subsection*{Bài 5}
	
	Một vật nhỏ có khối lượng $m$ được buộc vào đầu một sợi chỉ, đầu kia của sợi chỉ được luồn qua một lỗ nhỏ $O$ trên mặt bàn nhẵn nằm ngang, kéo căng sợi chỉ và giữ chặt đầu sợi chỉ. Khoảng cách từ $O$ tới vật là $r_0$. Truyền cho vật vận tốc vuông góc với sợi chỉ, vật chuyển động tròn đều trên mặt bàn với vận tốc $v_0$. Tại $t_0 = 0$, kéo nhẹ đều sợi chỉ qua lỗ. Sau thời gian $\tau$ sợi chỉ quay được góc $2\pi$ quanh $O$.
	
	% Để chèn hình ảnh, em cần có file ảnh (ví dụ: Bai9.png) 
	% trong cùng thư mục và bỏ dấu % ở dòng \usepackage{graphicx}
	% và các dòng bên dưới:
	
	% \begin{figure}[h!]
		%     \centering
		%     % \includegraphics[width=0.4\textwidth]{Bai9.png} 
		%     % Thay Bai9.png bằng tên file hình của em
		%     \caption{Hình minh họa Bài 9}
		%     \label{fig:bai9}
		% \end{figure}
	
	\begin{enumerate}
		\item Tìm diện tích quét được của sợi chỉ trong thời gian $\tau$.
		\item Tìm biểu thức lực căng của sợi chỉ theo khoảng cách $r$ từ $O$ tới vật.
		\item Tìm tỉ số giữa động năng của vật và tốc độ góc $\omega$ của sợi chỉ.
	\end{enumerate}
\begin{figure}[h!]
	\centering
	\includegraphics[width=0.5\linewidth]{"../../../itrithuc/OneDrive - Trong/OneDrive/0 AI LATEX/THU VIEN ANH/screenshot009"}
	\caption{}
	\label{fig:screenshot009}
\end{figure}

		% Giải thích các ký hiệu LaTeX:
	% $m$ : Biến số m (khối lượng) trong chế độ toán học
	% $r$ : Biến số r (khoảng cách)
	% $r = r_0$ : Phương trình với chỉ số dưới (subscript)
	% $\omega_0$ : Ký hiệu omega (\omega) với chỉ số dưới (_0)
	\vspace{1cm} % Tạo một khoảng trống
	
	% --- PHẦN BÀI GIẢI ---
\newpage
	\section*{Hướng dẫn giải chi tiết bài 3}
	\begin{figure}[ht]
		\centering
		\includegraphics[width=0.7\linewidth]{"C:/Users/admin/itrithuc/OneDrive - Trong/OneDrive/0 AI LATEX/THU VIEN ANH/screenshot001"}
		\caption{}
		\label{fig:screenshot001}
	\end{figure}
	
	\subsection*{Phân tích bài toán}
	Chúng ta có các thông tin sau:
	\begin{enumerate}
		\item \textbf{Hệ vật:} Gồm một vật nhỏ khối lượng $m$.
		\item \textbf{Môi trường:} Chuyển động trên mặt phẳng ngang, \textbf{nhẵn} (không có ma sát).
		\item \textbf{Lực tác dụng (theo phương ngang):} Chỉ có lực căng dây $T$ hướng vào tâm O.
		\item \textbf{Điều kiện:}
		\begin{itemize}
			\item Dây được kéo qua lỗ O với \textbf{vận tốc không đổi}. Điều này có nghĩa là vận tốc theo phương bán kính $v_r = \dot{r}$ là một hằng số. Vì dây được "kéo" (thu ngắn lại) nên $v_r$ là một giá trị âm không đổi.
			\item Từ $v_r = \dot{r} = \text{hằng số}$, ta suy ra gia tốc theo phương bán kính $\ddot{r} = 0$.
			\item Điều kiện ban đầu: Tại $r = r_0$, vận tốc góc là $\omega = \omega_0$.
		\end{itemize}
		\item \textbf{Yêu cầu:} Tìm sức căng $T$ của dây theo $r$.
	\end{enumerate}
	
	\subsection*{Hướng dẫn giải}
	Chúng ta sẽ sử dụng hai nguyên lý cơ bản:
	\begin{enumerate}
		\item Định luật II Newton cho chuyển động theo phương bán kính.
		\item Định luật Bảo toàn Mô-men Động lượng.
	\end{enumerate}
	
	\subsubsection*{Bước 1: Áp dụng Định luật Bảo toàn Mô-men Động lượng}
	Vì mặt phẳng nhẵn và lực căng $T$ luôn hướng về tâm O, mô-men lực tác dụng lên vật $m$ đối với tâm O bằng không:
	$$ \vec{\tau} = \vec{r} \times \vec{T} = 0 $$
	Do đó, mô-men động lượng $L$ của vật đối với tâm O được bảo toàn:
	$$ L =  (m r^2) \omega = \text{hằng số} $$
	Sử dụng điều kiện ban đầu ($r = r_0, \omega = \omega_0$) để tìm hằng số này:
	$$ L = m r_0^2 \omega_0 $$
	Vì $L$ bảo toàn tại mọi vị trí $r$, ta có:
	$$ m r^2 \omega = m r_0^2 \omega_0 $$
	Từ đây, ta tìm được biểu thức của vận tốc góc $\omega$ theo bán kính $r$:
	$$ \omega = \omega_0 \left( \frac{r_0}{r} \right)^2 = \frac{\omega_0 r_0^2}{r^2} $$
	
	\subsubsection*{Bước 2: Áp dụng Định luật II Newton (Phương bán kính)}
	Định luật II Newton cho chuyển động của vật trong hệ tọa độ cực (theo phương bán kính) là:
	$$ \sum F_r = m a_r = m (\ddot{r} - r\omega^2) $$
	\begin{itemize}
		\item Lực tác dụng duy nhất theo phương bán kính là lực căng dây $T$, hướng vào tâm (ngược chiều dương của $r$). Do đó, $\sum F_r = -T$.
		\item Vì dây được kéo với \textbf{vận tốc không đổi} nên gia tốc theo phương bán kính $\ddot{r} = 0$.
	\end{itemize}
	Thay các giá trị này vào phương trình:
	$$ -T = m (0 - r\omega^2) $$
	$$ T = m r \omega^2 $$
	Phương trình này cho thấy lực căng $T$ chính là lực hướng tâm.
	
	\subsubsection*{Bước 3: Kết hợp kết quả}
	Bây giờ, chúng ta thay biểu thức của $\omega$ (tìm được ở Bước 1) vào biểu thức của $T$ (tìm được ở Bước 2):
	$$ T = m r \left( \frac{\omega_0 r_0^2}{r^2} \right)^2 $$
	$$ T = m r \left( \frac{\omega_0^2 r_0^4}{r^4} \right) $$
	Rút gọn $r$, ta được kết quả cuối cùng:
	$$ T = m \omega_0^2 \frac{r_0^4}{r^3} $$
	
	\subsection*{Kết luận}
	Vậy, sức căng của dây theo khoảng cách $r$ được xác định bởi biểu thức:
	% \boxed{} được dùng để đóng khung công thức cuối cùng
	% Giải thích các ký hiệu LaTeX:
% $m$ : Biến số m (khối lượng) trong chế độ toán học
% $r$ : Biến số r (khoảng cách)
% $r = r_0$ : Phương trình với chỉ số dưới (subscript)
% $\omega_0$ : Ký hiệu omega (\omega) với chỉ số dưới (_0)
\vspace{1cm} % Tạo một khoảng trống

% --- PHẦN BÀI GIẢI ---
\section*{Hướng dẫn giải chi tiết bài 3}
\begin{figure}[ht]
	\centering
	\includegraphics[width=0.7\linewidth]{"C:/Users/admin/itrithuc/OneDrive - Trong/OneDrive/0 AI LATEX/THU VIEN ANH/screenshot001"}
	\caption{}
	\label{fig:screenshot001}
\end{figure}

\subsection*{Phân tích bài toán}
Chúng ta có các thông tin sau:
\begin{enumerate}
	\item \textbf{Hệ vật:} Gồm một vật nhỏ khối lượng $m$.
	\item \textbf{Môi trường:} Chuyển động trên mặt phẳng ngang, \textbf{nhẵn} (không có ma sát).
	\item \textbf{Lực tác dụng (theo phương ngang):} Chỉ có lực căng dây $T$ hướng vào tâm O.
	\item \textbf{Điều kiện:}
	\begin{itemize}
		\item Dây được kéo qua lỗ O với \textbf{vận tốc không đổi}. Điều này có nghĩa là vận tốc theo phương bán kính $v_r = \dot{r}$ là một hằng số. Vì dây được "kéo" (thu ngắn lại) nên $v_r$ là một giá trị âm không đổi.
		\item Từ $v_r = \dot{r} = \text{hằng số}$, ta suy ra gia tốc theo phương bán kính $\ddot{r} = 0$.
		\item Điều kiện ban đầu: Tại $r = r_0$, vận tốc góc là $\omega = \omega_0$.
	\end{itemize}
	\item \textbf{Yêu cầu:} Tìm sức căng $T$ của dây theo $r$.
\end{enumerate}

\subsection*{Hướng dẫn giải}
Chúng ta sẽ sử dụng hai nguyên lý cơ bản:
\begin{enumerate}
	\item Định luật II Newton cho chuyển động theo phương bán kính.
	\item Định luật Bảo toàn Mô-men Động lượng.
\end{enumerate}

\subsubsection*{Bước 1: Áp dụng Định luật Bảo toàn Mô-men Động lượng}
Vì mặt phẳng nhẵn và lực căng $T$ luôn hướng về tâm O, mô-men lực tác dụng lên vật $m$ đối với tâm O bằng không:
$$ \vec{\tau} = \vec{r} \times \vec{T} = 0 $$
Do đó, mô-men động lượng $L$ của vật đối với tâm O được bảo toàn:
$$ L = I \omega = (m r^2) \omega = \text{hằng số} $$
Sử dụng điều kiện ban đầu ($r = r_0, \omega = \omega_0$) để tìm hằng số này:
$$ L = m r_0^2 \omega_0 $$
Vì $L$ bảo toàn tại mọi vị trí $r$, ta có:
$$ m r^2 \omega = m r_0^2 \omega_0 $$
Từ đây, ta tìm được biểu thức của vận tốc góc $\omega$ theo bán kính $r$:
$$ \omega = \omega_0 \left( \frac{r_0}{r} \right)^2 = \frac{\omega_0 r_0^2}{r^2} $$

\subsubsection*{Bước 2: Áp dụng Định luật II Newton (Phương bán kính)}
Định luật II Newton cho chuyển động của vật trong hệ tọa độ cực (theo phương bán kính) là:
$$ \sum F_r = m a_r = m (\ddot{r} - r\omega^2) $$
\begin{itemize}
	\item Lực tác dụng duy nhất theo phương bán kính là lực căng dây $T$, hướng vào tâm (ngược chiều dương của $r$). Do đó, $\sum F_r = -T$.
	\item Vì dây được kéo với \textbf{vận tốc không đổi} nên gia tốc theo phương bán kính $\ddot{r} = 0$.
\end{itemize}
Thay các giá trị này vào phương trình:
$$ -T = m (0 - r\omega^2) $$
$$ T = m r \omega^2 $$
Phương trình này cho thấy lực căng $T$ chính là lực hướng tâm.

\subsubsection*{Bước 3: Kết hợp kết quả}
Bây giờ, chúng ta thay biểu thức của $\omega$ (tìm được ở Bước 1) vào biểu thức của $T$ (tìm được ở Bước 2):
$$ T = m r \left( \frac{\omega_0 r_0^2}{r^2} \right)^2 $$
$$ T = m r \left( \frac{\omega_0^2 r_0^4}{r^4} \right) $$
Rút gọn $r$, ta được kết quả cuối cùng:
$$ T = m \omega_0^2 \frac{r_0^4}{r^3} $$

\subsection*{Kết luận}
Vậy, sức căng của dây theo khoảng cách $r$ được xác định bởi biểu thức:
% \boxed{} được dùng để đóng khung công thức cuối cùng
	% Giải thích các ký hiệu LaTeX:
% $m$ : Biến số m (khối lượng) trong chế độ toán học
% $r$ : Biến số r (khoảng cách)
% $r = r_0$ : Phương trình với chỉ số dưới (subscript)
% $\omega_0$ : Ký hiệu omega (\omega) với chỉ số dưới (_0)
\vspace{1cm} % Tạo một khoảng trống

% --- PHẦN BÀI GIẢI ---
\section*{Hướng dẫn giải chi tiết bài 3}
\begin{figure}[ht]
	\centering
	\includegraphics[width=0.7\linewidth]{"C:/Users/admin/itrithuc/OneDrive - Trong/OneDrive/0 AI LATEX/THU VIEN ANH/screenshot001"}
	\caption{}
	\label{fig:screenshot001}
\end{figure}

\subsection*{Phân tích bài toán}
Chúng ta có các thông tin sau:
\begin{enumerate}
	\item \textbf{Hệ vật:} Gồm một vật nhỏ khối lượng $m$.
	\item \textbf{Môi trường:} Chuyển động trên mặt phẳng ngang, \textbf{nhẵn} (không có ma sát).
	\item \textbf{Lực tác dụng (theo phương ngang):} Chỉ có lực căng dây $T$ hướng vào tâm O.
	\item \textbf{Điều kiện:}
	\begin{itemize}
		\item Dây được kéo qua lỗ O với \textbf{vận tốc không đổi}. Điều này có nghĩa là vận tốc theo phương bán kính $v_r = \dot{r}$ là một hằng số. Vì dây được "kéo" (thu ngắn lại) nên $v_r$ là một giá trị âm không đổi.
		\item Từ $v_r = \dot{r} = \text{hằng số}$, ta suy ra gia tốc theo phương bán kính $\ddot{r} = 0$.
		\item Điều kiện ban đầu: Tại $r = r_0$, vận tốc góc là $\omega = \omega_0$.
	\end{itemize}
	\item \textbf{Yêu cầu:} Tìm sức căng $T$ của dây theo $r$.
\end{enumerate}

\subsection*{Hướng dẫn giải}
Chúng ta sẽ sử dụng hai nguyên lý cơ bản:
\begin{enumerate}
	\item Định luật II Newton cho chuyển động theo phương bán kính.
	\item Định luật Bảo toàn Mô-men Động lượng.
\end{enumerate}

\subsubsection*{Bước 1: Áp dụng Định luật Bảo toàn Mô-men Động lượng}
Vì mặt phẳng nhẵn và lực căng $T$ luôn hướng về tâm O, mô-men lực tác dụng lên vật $m$ đối với tâm O bằng không:
$$ \vec{\tau} = \vec{r} \times \vec{T} = 0 $$
Do đó, mô-men động lượng $L$ của vật đối với tâm O được bảo toàn:
$$ L = I \omega = (m r^2) \omega = \text{hằng số} $$
Sử dụng điều kiện ban đầu ($r = r_0, \omega = \omega_0$) để tìm hằng số này:
$$ L = m r_0^2 \omega_0 $$
Vì $L$ bảo toàn tại mọi vị trí $r$, ta có:
$$ m r^2 \omega = m r_0^2 \omega_0 $$
Từ đây, ta tìm được biểu thức của vận tốc góc $\omega$ theo bán kính $r$:
$$ \omega = \omega_0 \left( \frac{r_0}{r} \right)^2 = \frac{\omega_0 r_0^2}{r^2} $$

\subsubsection*{Bước 2: Áp dụng Định luật II Newton (Phương bán kính)}
Định luật II Newton cho chuyển động của vật trong hệ tọa độ cực (theo phương bán kính) là:
$$ \sum F_r = m a_r = m (\ddot{r} - r\omega^2) $$
\begin{itemize}
	\item Lực tác dụng duy nhất theo phương bán kính là lực căng dây $T$, hướng vào tâm (ngược chiều dương của $r$). Do đó, $\sum F_r = -T$.
	\item Vì dây được kéo với \textbf{vận tốc không đổi} nên gia tốc theo phương bán kính $\ddot{r} = 0$.
\end{itemize}
Thay các giá trị này vào phương trình:
$$ -T = m (0 - r\omega^2) $$
$$ T = m r \omega^2 $$
Phương trình này cho thấy lực căng $T$ chính là lực hướng tâm.

\subsubsection*{Bước 3: Kết hợp kết quả}
Bây giờ, chúng ta thay biểu thức của $\omega$ (tìm được ở Bước 1) vào biểu thức của $T$ (tìm được ở Bước 2):
$$ T = m r \left( \frac{\omega_0 r_0^2}{r^2} \right)^2 $$
$$ T = m r \left( \frac{\omega_0^2 r_0^4}{r^4} \right) $$
Rút gọn $r$, ta được kết quả cuối cùng:
$$ T = m \omega_0^2 \frac{r_0^4}{r^3} $$

\subsection*{Kết luận}
Vậy, sức căng của dây theo khoảng cách $r$ được xác định bởi biểu thức:
% \boxed{} được dùng để đóng khung công thức cuối cùng
$$ \mathbf{T = m \omega_0^2 \frac{r_0^4}{r^3}} $$
\end{document}